%==========================================================================
\section{Dependence on $h_\textup{ISCO}$, $\alpha$ and $q$}
\label{app:params}
%==========================================================================
% Figures expected in appendix/ (copies of plots_paper/pif_hr_alpha_map.pdf,
% produced by pif_hr_alpha.py, and of plots_paper/pif_hr_alpha_map_q05.pdf,
% produced by pif_q.py). All numbers from pif_hr_alpha.py and
% pif_constraints.py (spin step 0.005, r_out <= 200), and from pif_q.py for
% q = 1/2 (r_out <= 1000).

All the results of Sects.~\ref{sec:results} and \ref{sec:general} were obtained for $h_\textup{ISCO}=0.1$, $\alpha=0.01$ and $q=3/2$.
These are the values of \citet{Motta2017}, who adopted them for XRBs in the hard state.
% Check against Sect. 4.1 of Motta et al.: q = 3/2 is taken there as the profile of a radiation-pressure
% dominated disc, and (H/R) = 0.1 as appropriate to the hard state (accretion rate of 1-10% of Eddington);
% no justification is given for alpha = 0.01.
None of them depends explicitly on the mass.
In the radiation-pressure dominated region of a thin disc, for instance, the thickness is $H=\tfrac32\,(\lambda/\eta)\left[1-(r_\textup{in}/r)^{1/2}\right]R_g$ \citep{ss,accretion}, where $\lambda$ is the luminosity in units of the Eddington value and $\eta$ the radiative efficiency.
The aspect ratio is then a function of $r$, $\lambda$ and the spin only; its maximum, reached at $r=\tfrac94\,r_\textup{in}$, is $\tfrac29\,(\lambda/\eta)/r_\textup{in}\simeq0.6\,\lambda$ for $r_\textup{in}=r_\textup{ISCO}$ and both $a=0$ and $a=0.998$, so that $h_\textup{ISCO}=0.1$ corresponds to $\lambda\simeq0.15$--$0.2$.
The disc parameters of XRBs and AGNs are thus expected to differ where the accretion states differ, not because of the mass; they are in any case not constrained by observations in either class.
The equivalence of Sect.~\ref{sec2:scaling} between an XRB and an AGN with the same $P$ holds at fixed disc parameters, so that different parameters in the two classes are one way of breaking it.
In this appendix we therefore repeat the analysis of Sect.~\ref{sec:general} for other values of $h_\textup{ISCO}$ and $\alpha$ (Sects.~\ref{appA:scalings}--\ref{appA:sources}) and for a different thickness profile (Sect.~\ref{appA:q}), and examine how the constraints on the individual sources change.

%--------------------------------------------------------------------------
\subsection{Scalings}
\label{appA:scalings}

The two parameters enter the model only through constant factors.
The precession frequency, Eq.~\eqref{eq2:gdef}, does not depend on them.
By Eq.~\eqref{eq2:twave}, $\tau_w\propto h_\textup{ISCO}^{-1}$; by Eqs.~\eqref{eq2:talign} and \eqref{eq2:cs}, $\tau_a\propto c_s^2/\alpha\propto h_\textup{ISCO}^{2}/\alpha$.
The two sides of Eq.~\eqref{eq2:Ncond} thus scale as
\begin{equation}
    \label{eqA:scal}
    g\,\tau_w\propto h_\textup{ISCO}^{-1},
    \qquad
    g\,\tau_a\propto\frac{h_\textup{ISCO}^{2}}{\alpha},
\end{equation}
at every $(a,r_\textup{out})$.
Three consequences follow without any computation.
First, the upper edge of the valid region, $P<g_\textup{max}(a)$, is the same for all disc parameters, because Eq.~\eqref{eq2:Pcond} is, and because both conditions hold for $\bar r\to r_\textup{ISCO}$ whatever the constant factors.
Second, wave propagation depends on $h_\textup{ISCO}$ only: a thinner flow has a lower sound speed and is crossed more slowly.
Third, alignment depends on the single combination $h_\textup{ISCO}^{2}/\alpha$: a thicker or less viscous flow aligns more slowly.

We consider the grid $h_\textup{ISCO}\in\{0.03,0.1,0.3\}$, $\alpha\in\{10^{-3},10^{-2},10^{-1}\}$, over which $h_\textup{ISCO}^{2}/\alpha$ spans four orders of magnitude, from $9\times10^{-3}$ to 90, with the fiducial case in the middle.
By Eq.~\eqref{eqA:scal} this requires no new integration.
Figure~\ref{fig:hr_alpha_map} shows the $(a,P)$ plane of Fig.~\ref{fig:Pspin} for the nine combinations.

\begin{figure*}
    \centering
    \includegraphics[width=0.9\hsize]{appendix/pif_hr_alpha_map.pdf}
    \caption{
    As Fig.~\ref{fig:Pspin}, for $h_\textup{ISCO}\in\{0.03,0.1,0.3\}$ (columns) and $\alpha\in\{10^{-3},10^{-2},10^{-1}\}$ (rows); the fiducial case has a thick frame.
    The orange line is the value of $P$ for which $\bar r=r_\textup{ISCO}\,h_\textup{ISCO}/\alpha$, the outer limit of the bending-wave regime: configurations below it lie outside the regime.
    For $\alpha\geq h_\textup{ISCO}$ the line coincides with $g_\textup{max}(a)$; for $h_\textup{ISCO}/\alpha\geq100$ it lies outside the panel or beyond $r_\textup{out}=200$.
    }
    \label{fig:hr_alpha_map}
\end{figure*}

The validity of each panel is limited by the condition $\alpha\lesssim H/R$, i.e.\ $\bar r\lesssim r_\textup{ISCO}\,h_\textup{ISCO}/\alpha$ (Sect.~\ref{sec2:timescales}), which also depends on the two parameters.
The limit is $100$ and $300\,r_\textup{ISCO}$ for $\alpha=10^{-3}$ and $h_\textup{ISCO}\geq0.1$, and $30\,r_\textup{ISCO}$ for $(0.03,10^{-3})$ and $(0.3,10^{-2})$: in these four cases, as in the fiducial one, the valid region lies within the bending-wave regime, except at most for $|a|\lesssim0.1$.
For $(0.03,10^{-2})$ and $(0.3,10^{-1})$ the limit is $3\,r_\textup{ISCO}$, and only part of the valid region is within it: $|a|\gtrsim0.3$ in the first case, $a\gtrsim0.93$ in the second.
For $(0.1,10^{-1})$ and $(0.03,10^{-1})$, $\alpha\geq h_\textup{ISCO}$ and no radius is in the bending-wave regime: the warp would diffuse instead of propagating, Eq.~\eqref{eq2:talign} does not apply, and the corresponding panels are only formal.
The bottom row of Fig.~\ref{fig:hr_alpha_map} must therefore be read with caution.

%--------------------------------------------------------------------------
\subsection{The valid region}
\label{appA:general}

For all nine combinations $g$ is the same function, and the radii that satisfy Eq.~\eqref{eq2:Ncond} still form a single interval starting at the ISCO.
The description of Sect.~\ref{sec4:rmax} therefore holds unchanged: at fixed spin the valid configurations are $r_\textup{ISCO}<\bar r<r_\textup{max}(a)$, i.e.\ $P_\textup{min}(a)<P<g_\textup{max}(a)$, and the disc parameters only move $r_\textup{max}$, hence the lower edge $P_\textup{min}$.
In Fig.~\ref{fig:hr_alpha_map} the upper edge of the green region is the same in all panels, as are the cusp at $a=0$ and the limits $P<110$~Hz\,$M_\odot$ for retrograde and $P<1.62\times10^4$~Hz\,$M_\odot$ for prograde spins.
The grey band of Fig.~\ref{fig:rmax} is the range spanned by $r_\textup{max}$.

\paragraph{Alignment.}
Where alignment sets $r_\textup{max}$, the lower edge depends only on $h_\textup{ISCO}^{2}/\alpha$ (Table~\ref{tab:hr_alpha}).
Combinations with the same $h_\textup{ISCO}^{2}/\alpha$ within $10\%$ have the same $r_\textup{max}$ within about $5\%$: in Fig.~\ref{fig:hr_alpha_map} the panels along the diagonals from top left to bottom right differ only at high prograde spin, where wave propagation takes over.
For each decade in $h_\textup{ISCO}^{2}/\alpha$, $r_\textup{max}$ increases by a factor of $1.5$--$2.6$ and $P_\textup{min}$ decreases by $0.3$--$1$~dex, the more so the larger $h_\textup{ISCO}^{2}/\alpha$.
At $a=0.5$ the valid interval of $P$ is 0.4~dex wide for $h_\textup{ISCO}^{2}/\alpha=9\times10^{-3}$, 1.3~dex for the fiducial case and 3.1~dex for $h_\textup{ISCO}^{2}/\alpha=90$, and the largest valid flow extends to 7, 20 and 170~$R_g$, respectively.
A thin and viscous flow can thus precess visibly only if it is a narrow ring close to the ISCO, $\bar r\lesssim2\,r_\textup{ISCO}$, whereas a thick flow of low viscosity can extend to tens of $r_\textup{ISCO}$.
For retrograde spins the effect is the same: the lowest $P$ valid for $a=-1$ goes from 39 to less than 0.3~Hz\,$M_\odot$.

\begin{table}
\caption{Largest valid precession radius and lower limit of the valid interval of $P$ where they are set by alignment, as functions of $h_\textup{ISCO}^{2}/\alpha$.}
\label{tab:hr_alpha}
\centering
\begin{tabular}{lcccc}
\hline\hline
$h_\textup{ISCO}^{2}/\alpha$ & \multicolumn{2}{c}{$a=-1$} & \multicolumn{2}{c}{$a=0.5$} \\
 & $r_\textup{max}$ & $P_\textup{min}$ & $r_\textup{max}$ & $P_\textup{min}$ \\
\hline
$9\times10^{-3}$ & 15       & 39          & 7.1      & 137 \\
$0.09$--$0.1$    & 23       & 18          & 10--11   & 64--67 \\
$0.9$--$1$       & 43--45   & 4.9--5.2    & 20       & 19--20 \\
$9$--$10$        & 113--119 & 0.72--0.80  & 50--53   & 2.9--3.2 \\
$90$             & $>200$   & $<0.26$     & 171      & 0.29 \\
\hline
\end{tabular}
\tablefoot{
Radii in $R_g$, $P_\textup{min}$ in Hz\,$M_\odot$; for comparison $g_\textup{max}=110$ and $347$~Hz\,$M_\odot$ for $a=-1$ and $0.5$.
Rows with a range contain two or three combinations of the grid: $(h_\textup{ISCO},\alpha)=(0.03,10^{-2})$ and $(0.1,10^{-1})$; $(0.1,10^{-2})$ and $(0.3,10^{-1})$; $(0.1,10^{-3})$ and $(0.3,10^{-2})$.
The combination $(0.03,10^{-3})$, with $h_\textup{ISCO}^{2}/\alpha=0.9$, is omitted at $a=0.5$, where it is limited by wave propagation ($r_\textup{max}=18$, $P_\textup{min}=24$~Hz\,$M_\odot$).
}
\end{table}

\paragraph{Wave propagation.}
Wave propagation limits the valid region only for rapidly spinning prograde holes, from a spin $a_w$ upwards (dark grey in Fig.~\ref{fig:hr_alpha_map}).
Since it becomes the limiting condition where it fails before alignment, $a_w$ decreases with $h_\textup{ISCO}$ and with $\alpha$: $a_w=0.48$, $0.65$ and $0.80$ for $h_\textup{ISCO}=0.03$ and $\alpha=10^{-3}$, $10^{-2}$, $10^{-1}$; $a_w=0.77$, $0.91$ and $0.965$ for $h_\textup{ISCO}=0.1$; $a_w=0.95$ and $0.998$ for $h_\textup{ISCO}=0.3$ and $\alpha=10^{-3}$, $10^{-2}$, while for $(0.3,10^{-1})$ alignment is the limit at all spins.
Above $a_w$ the lower edge depends on $h_\textup{ISCO}$ only.
For $h_\textup{ISCO}=0.03$ the valid flows around rapidly spinning holes are very narrow rings, $r_\textup{max}=1.24\,r_\textup{ISCO}$ at $a=0.9$ and $1.03\,r_\textup{ISCO}$ at $a=0.998$, and the valid interval of $P$ is only 0.15 and 0.02~dex wide.
For $h_\textup{ISCO}=0.3$, instead, wave propagation is almost irrelevant: at $a=0.998$ the flow can extend to $2.3$--$3.9\,r_\textup{ISCO}$, and the valid interval is $0.6$--$1$~dex wide.
The narrowing of the valid region towards $a\to1$ found in Sect.~\ref{sec4:aP} is therefore a property of thin flows.

\paragraph{Valid spins at fixed $P$.}
In terms of the set $A(P)$ of Sect.~\ref{sec4:AP}, the end of each interval closer to $a=0$, which is set by $g_\textup{max}(a)=P$, does not move, while the other end does.
The minimum $|a|$ required by a given $P$, and the absence of retrograde solutions above $110$~Hz\,$M_\odot$, are thus independent of the disc parameters.
The maximum valid spin increases with $h_\textup{ISCO}^{2}/\alpha$ and with $h_\textup{ISCO}$.
At $P=100$~Hz\,$M_\odot$, where the minimum prograde spin is $0.24$, it goes from $0.44$ for $(0.03,10^{-1})$ to $0.76$ for the fiducial case and $0.98$ for $(0.3,10^{-3})$.
At $P=10^3$~Hz\,$M_\odot$, where the minimum is $0.75$, it depends almost only on the thickness: $0.84$ for $h_\textup{ISCO}=0.03$, $0.92$--$0.94$ for $0.1$, and $0.97$--$0.995$ for $0.3$.
At low $P$ the change is the largest.
At $P=1$~Hz\,$M_\odot$ the valid spins are $|a|\lesssim0.03$ for $(0.03,10^{-1})$, $a\leq0.16$ and $a\geq-0.23$ for the fiducial case, and $a\leq0.35$--$0.66$ with all retrograde spins for $h_\textup{ISCO}^{2}/\alpha\geq9$.
The correspondence between $P$ and spin of Sect.~\ref{sec4:AP} is therefore robust at high $P$, where a rapidly spinning hole is required in all cases, and becomes weak for $P\lesssim10$~Hz\,$M_\odot$.

%--------------------------------------------------------------------------
\subsection{Individual sources}
\label{appA:sources}

Each source selects an interval of $P$ in Fig.~\ref{fig:hr_alpha_map}, so that the constraints of Table~\ref{tab:constraints} change as follows.

The lower bounds on $|a|$ and the upper bounds on the mass are the same for all nine combinations, since they follow from $g_\textup{max}$ (Sect.~\ref{sec3:grids}).
In particular, the need for a prograde spin in the ten sources of Sect.~\ref{sec3:highP}, their minimum spins, and the mass limits of NGC~4151, RE~J1034+396, MCG-06-30-15, 1H~0707-495, NGC~1365 and 1ES~1927+654 do not depend on the disc parameters.
% Numerical note: for h_ISCO = 0.03 the valid interval at a ~ 0.998 is narrower than the
% r_out grid spacing (2.3%), so the grid loses the last ~0.1 dex below the mass limit
% (e.g. the valid fraction of the mass range of RE J1034+396 drops from 0.91 to 0.85).
% The continuous limit is unchanged.

For the same ten sources the upper bound on the spin changes by about $0.1$, mostly with $h_\textup{ISCO}$.
Where the mass range does not reach $M_\textup{max}$, the maximum spin is $0.85$ for $h_\textup{ISCO}=0.03$ and $0.97$--$0.998$ for $h_\textup{ISCO}=0.3$ for Mrk~766 ($0.95$ in the fiducial case), $0.89$ and $0.98$--$0.998$ for MS~2254.9-3712 ($0.96$), and $0.86$ and $0.98$--$0.998$ for the higher-frequency QPO of ESO~113-G010 ($0.95$).
The precessing flow remains compact, $\bar r<15\,R_g$ in all cases.
A fast but not necessarily extremal spin and a compact flow are thus required for all combinations.

For the five sources of Sect.~\ref{sec3:lowP} the change is larger.
The maximum prograde spin varies between $0.57$ and $0.99$ for J1257, between $0.46$ and $0.98$ for the lower-frequency QPO of J1347, and between $0.65$ and $0.99$ for Mrk~142, while their retrograde solutions do not change.
The two sources with the lowest $P$ are the most sensitive.
For J1231 the maximum prograde spin goes from $0.11$ to $0.89$, and retrograde spins are limited to $|a|\leq0.14$--$0.36$ for $h_\textup{ISCO}^{2}/\alpha\leq0.1$.
For NGC~4945 the valid spins go from $|a|\leq0.04$ for $(0.03,10^{-1})$ to $a\leq0.74$ and all retrograde spins for $(0.3,10^{-3})$; retrograde spins down to $a=-1$ are valid only for $h_\textup{ISCO}^{2}/\alpha\geq9$.
The conclusion that these two sources require a slowly spinning hole (Sect.~\ref{sec3:lowP}) therefore holds only for $h_\textup{ISCO}^{2}/\alpha\lesssim1$.

The source most affected is 1ES~1927+654 (Sect.~\ref{sec3:pairs}).
Its two frequencies differ by a factor of $2.7$, i.e.\ $0.43$~dex in $P$ at fixed mass, and both require $a\geq0.84$.
A common spin exists only if the valid interval of $P$ at that spin is wider than $0.43$~dex.
For $h_\textup{ISCO}=0.03$ it is less than $0.3$~dex wide for $a\geq0.84$, and no spin reproduces both frequencies, whatever $\alpha$: in the model, the observed drift excludes a thin flow.
For $h_\textup{ISCO}=0.1$ the solution is that of Eq.~\eqref{eq3:1es}, with the mass limit reduced to $1.3\times10^6\,M_\odot$ for $\alpha=10^{-1}$.
For $h_\textup{ISCO}=0.3$ it widens to $0.85\lesssim a\leq0.998$ and $M\lesssim4.1\times10^6\,M_\odot$.
The joint solutions of the simultaneous pairs survive in all cases, but shrink for thin flows: for 1H~0707-495, $0.62\lesssim a\lesssim0.77$ and $M\lesssim4\times10^6\,M_\odot$ for $h_\textup{ISCO}=0.03$ and $\alpha\leq10^{-2}$.

%--------------------------------------------------------------------------
\subsection{Thickness profile}
\label{appA:q}
% All numbers and Fig. hr_alpha_map_q from pif_q.py (q = 1/2).

The exponent $q$ of Eq.~\eqref{eq2:HR} does not enter the model through a constant factor.
It sets the radial dependence of the sound speed (Eq.~\ref{eq2:cs}), hence the integrands of Eqs.~\eqref{eq2:twave} and \eqref{eq2:talign}, and its effect cannot be obtained from the fiducial results by a rescaling.
For $q=3/2$ the thickness does not depend on radius and the aspect ratio decreases outwards, $H/R=h_\textup{ISCO}\,r_\textup{ISCO}/r$: for $h_\textup{ISCO}=0.1$ it is below $10^{-2}$ beyond $10\,r_\textup{ISCO}$.
This is the profile of a radiation-pressure dominated thin disc, not that of a geometrically thick hot flow, whose aspect ratio is roughly constant \citep[$H/R\simeq0.2$ in][]{PIF}.
% Check that Ingram, Done & Fragile (2009) adopt h/r = 0.2 for the hot flow.
We therefore repeat the analysis for $q=1/2$, i.e.\ $H/R=h_\textup{ISCO}$ at every radius.

Three properties follow without any computation.
First, the precession frequency does not depend on $q$: $g(a,r_\textup{out})$, the upper edge $g_\textup{max}(a)$, the minimum spins and the maximum masses are unchanged, and Eq.~\eqref{eqA:scal} still holds exactly.
Second, since $(r/r_\textup{ISCO})^{-1/2}\geq(r/r_\textup{ISCO})^{-3/2}$ for $r\geq r_\textup{ISCO}$, the sound speed is larger at every radius than for $q=3/2$ and the same $h_\textup{ISCO}$.
Then $t_\textup{wave}$ is shorter and $t_\textup{align}$ is longer (in Eq.~\ref{eq2:talign}, $\mathcal{G}$ and $\Sigma$ do not depend on $c_s$), and every configuration that is valid for $q=3/2$ is valid for $q=1/2$.
Third, the condition $\alpha\lesssim H/R$ becomes $\alpha\lesssim h_\textup{ISCO}$ at all radii: the whole flow is in the bending-wave regime, or none of it is.

Since the valid region reaches larger radii, we extend the grid of $r_\textup{out}$ to $10^3\,R_g$ with the same logarithmic spacing, and for $r_\textup{out}>200$ we increase the number of radial nodes so that their spacing does not exceed $0.1\,R_g$, as in the fiducial grid.
With four times as many nodes, $r_\textup{max}$ and $P_\textup{min}$ change by less than $1\%$.
Figure~\ref{fig:hr_alpha_map_q} shows the $(a,P)$ plane for the nine combinations of Sect.~\ref{appA:scalings}, and Table~\ref{tab:q} lists $r_\textup{max}$ and $P_\textup{min}$.

\begin{figure*}
    \centering
    \includegraphics[width=0.9\hsize]{appendix/pif_hr_alpha_map_q05.pdf}
    \caption{
    As Fig.~\ref{fig:hr_alpha_map}, for a flow of constant aspect ratio, $q=1/2$ ($H/R=h_\textup{ISCO}$ at all radii).
    The upper edge of the valid region is the same as for $q=3/2$.
    Medium grey: the precession radius lies beyond $r_\textup{out}=10^3$.
    For $\alpha<h_\textup{ISCO}$ the whole flow is in the bending-wave regime; for $\alpha\geq h_\textup{ISCO}$ no radius is, and the orange line coincides with $g_\textup{max}(a)$.
    }
    \label{fig:hr_alpha_map_q}
\end{figure*}

\begin{table*}
\caption{Largest valid precession radius and lower limit of the valid interval of $P$ for a constant thickness ($q=3/2$) and for a constant aspect ratio ($q=1/2$).}
\label{tab:q}
\centering
\begin{tabular}{rcccccccc}
\hline\hline
 & \multicolumn{2}{c}{$q=3/2$} & \multicolumn{6}{c}{$q=1/2$} \\
 & \multicolumn{2}{c}{$(0.1,\,10^{-2})$} & \multicolumn{2}{c}{$(0.1,\,10^{-2})$} & \multicolumn{2}{c}{$(0.2,\,10^{-2})$} & \multicolumn{2}{c}{$(0.3,\,10^{-1})$} \\
$a$ & $r_\textup{max}$ & $P_\textup{min}$ & $r_\textup{max}$ & $P_\textup{min}$ & $r_\textup{max}$ & $P_\textup{min}$ & $r_\textup{max}$ & $P_\textup{min}$ \\
\hline
$-1$    & 45   & 4.9              & $>10^3$ & $<0.012$          & $>10^3$ & $<0.012$  & 711     & 0.022 \\
$-0.5$  & 45   & 2.7              & $>10^3$ & $<0.007$          & $>10^3$ & $<0.007$  & $>10^3$ & $<0.007$ \\
$0.1$   & 59   & 0.39             & $>10^3$ & $<0.002$          & $>10^3$ & $<0.002$  & $>10^3$ & $<0.002$ \\
$0.3$   & 31   & 4.5              & $>10^3$ & $<0.006$          & $>10^3$ & $<0.006$  & $>10^3$ & $<0.006$ \\
$0.5$   & 20   & 19               & 316     & 0.090             & $>10^3$ & $<0.010$  & 215     & 0.19 \\
$0.7$   & 13   & 66               & 58      & 3.9               & $>10^3$ & $<0.016$  & 49      & 5.4 \\
$0.9$   & 7.4  & 328              & 17      & 66                & 408     & 0.15      & 16      & 79 \\
$0.95$  & 3.9  & $1.24\times10^3$ & 11      & 171               & 110     & 2.2       & 10.5    & 200 \\
$0.98$  & 2.3  & $3.39\times10^3$ & 7.8     & 408               & 48      & 13        & 7.2     & 467 \\
$0.998$ & 1.39 & $8.96\times10^3$ & 1.41    & $8.84\times10^3$  & 18      & 94        & 4.1     & $1.56\times10^3$ \\
\hline
\end{tabular}
\tablefoot{
The pairs in the header are $(h_\textup{ISCO},\alpha)$.
Radii in $R_g$, $P_\textup{min}$ in Hz\,$M_\odot$; $g_\textup{max}$ is given in Table~\ref{tab:window}.
Where $r_\textup{max}>10^3$ the valid region reaches the largest radius of the grid, and the upper limit on $P_\textup{min}$ is $g(a,10^3)$.
All the limits are set by alignment, except for $a\geq0.95$ with $q=3/2$ and for $a=0.998$ with $q=1/2$ and $h_\textup{ISCO}=0.1$, where they are set by wave propagation.
}
\end{table*}

\paragraph{Alignment.}
For $q=3/2$ the number of cycles $N_a=g\,\tau_a$ (Eq.~\ref{eq5:Na}) decreases steeply with radius, whereas for $q=1/2$ it flattens: for $a=0.5$ and the fiducial $h_\textup{ISCO}$ and $\alpha$, $N_a=35$, $2.3$, $1.1$ and $0.8$ at $r_\textup{out}=10$, $50$, $200$ and $10^3$, against $13$, $0.11$, $8.5\times10^{-3}$ and $6\times10^{-4}$.
Two consequences follow.
The first is that the valid region is much wider.
For the fiducial $h_\textup{ISCO}$ and $\alpha$, alignment does not limit the flow within $10^3\,R_g$ for all retrograde spins and for $a\leq0.43$; at $a=0.5$, $r_\textup{max}$ increases from 20 to $316\,R_g$ and the valid interval of $P$ from 1.3 to 3.6~dex.
For higher spins the change is smaller but still large: $r_\textup{max}$ increases by a factor of $2.3$--$4.4$ for $0.7\leq a\leq0.98$, and $P_\textup{min}$ decreases by $0.7$--$1.2$~dex.
The second is that, where it is large, $r_\textup{max}$ is very sensitive to $h_\textup{ISCO}^2/\alpha$, because $N_a$ crosses unity at a shallow angle.
A factor of four in $h_\textup{ISCO}^2/\alpha$ moves $r_\textup{max}$ at $a=0.9$ from 17 to $408\,R_g$, and at $a=0.5$ from 316 to beyond the grid.
For $q=1/2$ alignment is therefore better described as a condition on $h_\textup{ISCO}^2/\alpha$ at each spin than as a limit on the radius: for $h_\textup{ISCO}^2/\alpha\geq9$ all the spins up to $a=0.998$ are valid at $P=1$~Hz\,$M_\odot$, and for $h_\textup{ISCO}^2/\alpha=0.9$--$1$ those with $a\leq0.59$--$0.61$ (against $a\leq0.15$--$0.16$ for $q=3/2$).

\paragraph{Wave propagation.}
For $q=1/2$ the number of wave crossings per period, $g\,\tau_w$, is not monotonic: it has a maximum at $\bar r\simeq1.8$--$2.7\,r_\textup{ISCO}$ and decreases at larger radii, because the crossing time increases more slowly than the precession period.
Bending waves therefore cross the flow in time at all radii if they do so near the maximum, which requires $h_\textup{ISCO}>0.05$, $0.07$, $0.10$ and $0.19$ for $a=0.9$, $0.95$, $0.98$ and $0.998$.
For $h_\textup{ISCO}\geq0.2$ wave propagation never limits the valid region.
For $h_\textup{ISCO}=0.1$ it fails only for $a\geq0.985$, in an interval of radii around the maximum: $2.4\lesssim\bar r\lesssim4.4\,R_g$ for $a=0.985$ and $2.0\lesssim\bar r\lesssim5.3\,R_g$ for $a=0.99$.
For these two spins the valid radii form two intervals instead of one, so that the description of Sect.~\ref{sec4:rmax} does not apply; for $a\geq0.995$ alignment fails before wave propagation recovers, and the flow is limited to $\bar r<1.1$--$1.2\,r_\textup{ISCO}$ as for $q=3/2$.
For $h_\textup{ISCO}=0.03$ the same happens from $a\simeq0.8$--$0.85$ upwards (dark grey and white regions in the left column of Fig.~\ref{fig:hr_alpha_map_q}).

\paragraph{Rapidly spinning holes at low $P$.}
The region most affected is that of high spin and low $P$, which is the one occupied by the XRBs with a reflection spin $a>0.9$ and a type-C QPO (Sect.~\ref{sec5:caveats}).
At fixed $a$ and $P$ the precession radius does not depend on the disc parameters, and by Eq.~\eqref{eqA:scal} the conditions of Eq.~\eqref{eq2:Ncond} are equivalent to
\begin{equation}
    \label{eqA:bounds}
    h_\textup{ISCO}>0.1\,N_w^\textup{fid},
    \qquad
    \alpha<0.01\left(\frac{h_\textup{ISCO}}{0.1}\right)^{2}N_a^\textup{fid},
\end{equation}
where $N_w^\textup{fid}$ and $N_a^\textup{fid}$ are the values of $g\,\tau_w$ and $g\,\tau_a$ at $\bar r(a,P)$ for $h_\textup{ISCO}=0.1$ and $\alpha=0.01$.
Table~\ref{tab:xrb} lists the two limits for the two profiles.
For $q=3/2$ a hole with $a=0.9$ is valid at $P=100$~Hz\,$M_\odot$ for $h_\textup{ISCO}=0.2$--$0.3$ and $\alpha\lesssim0.006$--$0.014$, but at $P=10$~Hz\,$M_\odot$ it requires $\alpha\lesssim10^{-3}$, and at $P=1$~Hz\,$M_\odot$ it requires $h_\textup{ISCO}\simeq0.5$; for $a=0.998$ it would require $h_\textup{ISCO}\gtrsim0.5$--$1.7$, which is not a disc.
For $q=1/2$ wave propagation is satisfied for $h_\textup{ISCO}>0.06$ in all cases, and for $h_\textup{ISCO}=0.2$ alignment is satisfied for $\alpha\lesssim0.01$--$0.05$ at $a=0.9$ and $\alpha\lesssim0.005$--$0.01$ at $a=0.998$.
A viscosity $\alpha=0.1$ is instead excluded at these spins for $P\leq10$~Hz\,$M_\odot$ even for $h_\textup{ISCO}=0.3$, and is allowed at $P=100$~Hz\,$M_\odot$ only for $a\lesssim0.9$.

\begin{table}
\caption{Limits on $h_\textup{ISCO}$ and $\alpha$ for rapidly spinning holes at the values of $P$ of type-C QPOs.}
\label{tab:xrb}
\centering
\begin{tabular}{rrrccccc}
\hline\hline
 & & & \multicolumn{2}{c}{$q=3/2$} & \multicolumn{3}{c}{$q=1/2$} \\
$a$ & $P$ & $\bar r$ & $h_\textup{min}$ & $\alpha_\textup{max}$ & $h_\textup{min}$ & \multicolumn{2}{c}{$\alpha_\textup{max}$} \\
 & & & & $(0.3)$ & & $(0.2)$ & $(0.3)$ \\
\hline
$0.9$   & 1   & 151 & 0.48 & --                  & 0.012 & 0.012  & 0.028 \\
        & 10  & 46  & 0.25 & $1.1\times10^{-3}$  & 0.021 & 0.019  & 0.043 \\
        & 100 & 14  & 0.13 & $1.4\times10^{-2}$  & 0.036 & 0.051  & 0.12 \\
$0.95$  & 1   & 165 & 0.71 & --                  & 0.014 & 0.0090 & 0.020 \\
        & 10  & 50  & 0.37 & --                  & 0.024 & 0.013  & 0.030 \\
        & 100 & 15  & 0.20 & $5.7\times10^{-3}$  & 0.043 & 0.029  & 0.066 \\
$0.998$ & 1   & 195 & 1.7  & --                  & 0.018 & 0.0048 & 0.011 \\
        & 10  & 59  & 0.86 & --                  & 0.031 & 0.0062 & 0.014 \\
        & 100 & 18  & 0.46 & --                  & 0.057 & 0.010  & 0.023 \\
\hline
\end{tabular}
\tablefoot{
$P$ in Hz\,$M_\odot$ ($P=1$, 10 and 100 correspond to 0.1, 1 and 10~Hz for $M=10\,M_\odot$); $\bar r$ in $R_g$, independent of $q$.
$h_\textup{min}$: smallest $h_\textup{ISCO}$ for which bending waves cross the flow within one period.
$\alpha_\textup{max}$: largest $\alpha$ for which the flow does not align within one period, for the value of $h_\textup{ISCO}$ in parentheses (Eq.~\ref{eqA:bounds}); a dash means that this value is below $h_\textup{min}$.
}
\end{table}

\paragraph{Individual sources.}
The lower bounds on $|a|$ and the upper bounds on the mass do not change.
For the ten sources of Sect.~\ref{sec3:highP} the changes are small, because their flows are compact: for the fiducial $h_\textup{ISCO}$ and $\alpha$ the largest valid spin increases from $0.95$--$0.96$ to $0.99$ for Mrk~766, MS~2254.9-3712 and ESO~113-G010, and the largest precession radius by up to $\simeq20\%$, to at most $13\,R_g$.
A fast prograde spin and a compact flow are still required.
The joint solution of 1ES~1927+654 widens to $0.85\lesssim a\lesssim0.985$ and $M\lesssim2.4\times10^6\,M_\odot$, and to $a\leq0.998$ and $M\lesssim4.1\times10^6\,M_\odot$ for $h_\textup{ISCO}=0.2$.
The five sources of Sect.~\ref{sec3:lowP} change more.
The largest valid prograde spin increases from $0.84$ to $0.96$ for J1257, from $0.89$ to $0.96$ for Mrk~142 and from $0.78$ to $0.93$ for the lower-frequency QPO of J1347.
For J1231 it increases from $0.38$ to $0.75$, with precession radii up to $71\,R_g$, and for NGC~4945 from $0.21$ to $0.66$, with all the retrograde spins valid and precession radii up to $234\,R_g$; for $h_\textup{ISCO}=0.2$ and $\alpha=0.01$ both sources are valid up to $a\simeq0.95$--$0.97$.
The conclusion that these two sources require a slowly spinning hole (Sect.~\ref{sec3:lowP}) is therefore specific to a flow of constant thickness.

%--------------------------------------------------------------------------
\subsection{Summary}
\label{appA:summary}

The constraints that follow from the precession frequency alone, i.e.\ the upper edge of the valid region, the minimum spins and the maximum masses, do not depend on $h_\textup{ISCO}$, $\alpha$ and $q$.
The disc parameters set how far from the ISCO a valid flow can extend, hence the lowest valid $P$ at each spin: alignment through $h_\textup{ISCO}^{2}/\alpha$ at all spins, wave propagation through $h_\textup{ISCO}$ for rapidly spinning prograde holes.
Thin or viscous flows are valid only as narrow rings close to the ISCO, in an interval of $P$ of a few tenths of a dex below $g_\textup{max}(a)$; thick flows of low viscosity are valid over two to three decades in $P$.
The thickness profile has the largest effect.
If the aspect ratio is constant instead of decreasing outwards, the lower edge moves down by $0.7$--$1.2$~dex for $0.7\leq a\leq0.98$ and by more than two decades for $a\leq0.5$, wave propagation is no longer a constraint for $h_\textup{ISCO}\geq0.2$, and rapidly spinning holes are valid at the values of $P$ of type-C QPOs for $h_\textup{ISCO}=0.2$--$0.3$ and $\alpha\lesssim0.005$--$0.03$, depending on the spin.
The lower edge of the valid region, and the upper bounds on the spin that follow from it, are thus the results that depend most on the assumed structure of the flow.

With respect to the question of Sect.~\ref{sec:intro}, a difference in the disc parameters between XRBs and AGNs does not move the upper edge, and so it cannot make accessible to one class values of $P$ that are forbidden to the other.
It changes instead the extent of the valid region below the edge, i.e.\ the range of spins compatible with a given $P$ and the range of $P$ compatible with a given spin.
If the inner flows of AGNs were thinner or more viscous than those of XRBs, or if their thickness were constant while the aspect ratio of the flows of XRBs is, at a given spin a QPO could appear only in a narrow interval of $P$ just below $g_\textup{max}(a)$, and only for a flow confined to the immediate vicinity of the ISCO.
